
\chapter{Kontexterkenner}

In diesem Kapitel besch"aftige ich mit der Frage der Erkennung von Zusammenh"angen im \www und wie man diese erkennen kann.

Die hier vorgestellten Algorithmen wurden von mir nur in Einzelversuchen auf ihre G"ultigkeit "uberpr"uft. Da diese Algorithmen sehr speicherintensiv (sowohl Hauptspeicher als auch Plattenspeicher) sind, war es mir nicht m"oglich eine allgemein benutzbare Implementierung zu schaffen.


\section{"Uberlegungen}

Wie schon in der Einleitung erw"ahnt, bietet die gro"se Informationsmenge die das \www darstellt neue M"oglichkeiten der Informationsverabeitung. Kontexterkennung bezieht sich hierbei auf das weltweit verteilte Umfeld von Informationen.

Vor einigen Jahren publizierten Wissenschaftler neueste Erkenntnisse ihrer Forschung in Fachzeitschriften. Seit Einf"uhrung des \www verlagert sich die Publikation neuester Forschungsberichte zunehmend auf das Internet.

Da all diese Daten in aller Welt verstreut liegen und Suchserver momentan nur nach Stichworten suchen lassen, ist die Suche nach Hypertextdokumenten, die zu einem Themengebiet geh"oren, recht m"uhsam.

Abhilfe k"onnten hier jedoch die nun folgenden Strategien bieten.



\subsection{Domainerkennung und "`hot topics"'}

Mit dem im vorigen Kapitel vorgestellten Sprachenerkenner kann man eine Domainerkennung implementieren. 

Domain meint hier einen abgegrenzten Wortschatz. So kann man medizinische Texte von mathematischen Texten durch den verwendeten Wortschatz unterscheiden. Hier kann man das Themengebiete mit W"orterb"uchern beschreiben und Texte werden dann nicht mehr einer Sprache zugeh"orig angesehen, sondern einem Themengebiet.

Allerdings erleichtert dies einem die Suche nicht, sondern erm"oglicht es neue Seiten, die zu einem Themengebiet hinzugef"ugt werden, zu erkennen. In einem solchen Szenario k"onnte man sich vorstellen, da"s man so "`brandaktuelle"' Themen innerhalb eines Themenkomplexes erkennt.

F"ur ein noch nicht erw"ahntes "`Goodie"' der Spracherkennung kann man die Frequenzlisten aus \ref{WoerterbuecherUndWortfrequenzen} benutzen. Sehr wichtige W"orter sind hochfrequent, da h"aufig benutzt.

Wenn nun Worte in der Frequenzliste "`wandern"', d.h. in ihrer Wichtigkeit varieren, kann man erkennen, welche Worte akutelle Schlagworte sind. Man kann aber auch zuk"unftig Interessannte Themen schon vorher erkennen, da diese Worte in Fachkreisen schon vorher h"aufiger benutzt werden.

Benutzt man nun einen {\em Domainerkenner} und verfolgt dort das "`Wandern"' von Worten, so kann man neue Entwicklungen in den verschiedenen Themengebieten von Anfang an mitverfolgen.



\subsection{Clustering Algorithmen}

Diese Strategie wird momentan vom {\em AltaVista-LiveTopics} Suchdienst benutzt. Man sucht nach einem Wort und es wird in einigen tausend Dokumenten gefunden. Diese Dokumente werden dann durch Clustering Algorithmen in zusammenh"angende Gruppen unterteilt.

Dies funktioniert einfach ausgedr"uckt so:

Aus allen Dokumenten die das Suchwort enthalten werden jeweils alle Worte herausgefiltert. Dann werden die Worte, welche auf mehren Seiten zusammen auftauchen, in Gruppen zusammengefa"st. Die "Uberg"ange zwischen den Gruppen werden dann als verschiedene Bereiche des Suchbegriffes aufgefa"st.

Da man alle Hypertextdokumente mit allen anderen Hypertextdokumenten des Ergebnisraums vergleichen mu"s ( z.B. 172.000 Dokumente beim Wort "`composer"' ), ben"otigt man bei Durchf"uhrung dieser Methode sehr viel Rechen- und Speicherkapazit"at.


\section{Normalisierung}

Normalisieren (lemmatisieren\footnote{Lemmatisieren wird das Zerlegen von Worten in die einzelnen Komponenten genannt.}) von Hypertextdokumenten bedeutet, da"s man die Worte in den Dokumenten in die jeweiligen Grundformen zerlegt. Erst diese Grundformen benutzt man dann zum indexieren.

An einem Beispiel wird dies ersichtlicher: Wenn man den das Wort "'Komponisten"' indexiert, dann kann man dieses Dokument nur "uber Eingabe des Wortes "`Komponisten"', nicht jedoch "uber Eingabe von "`Komponist"' erreichen.

Normalisiert man das Wort "`Komponisten"' zu "`Komponist"' und normalisiert dar"uber hinaus noch die Eingaben der Suchanfragen, so bekommt man bei Suchanfragen auch Seiten, die das Suchwort nicht direkt enthalten.

Allerdings setzt dies auch voraus, da"s man zum einen einen funktionierenden Sprachenerkenner hat und zum anderen einen Normalisierer f"ur die entsprechende Sprache. Am {\em CIS} wurde ein solcher Normalisierer f"ur die deutsche Sprache schon implementiert, so da"s ich darauf schon zurr"uckgreifen konnte.

Diese Normalisierung bietet noch einen Vorteil bei der Bearbeitung von Hypertextdokumenten an: Die Indexe werden wesentlich kleiner, da weniger verschiedene Worte zu indexieren sind.


\subsection{Linkauswertung}

Weniger Rechenleistung ben"otigt man, wenn man sich die Intelligenz der Leute im Netz zunutze machen kann. So kennen die Authoren von Hypertextdokumenten schon einen kleinen Teil des \www und referenzieren diesen. Die Linkauswertung verlangt insofern wesentlich geringere Rechenkapazit"aten, als da"s die Links in der Suchmaschine sowieso vorgehalten werden m"ussen und nur die Bezeichner von Verweisen zus"atzlich gespeichert werden m"ussen.

Man findet auf Seiten zu "`Mozart"' Referenzen auf andere "`Mozart"'-Seiten, Seiten zu anderen klassischen und modernen Komponisten und viele in das Themengebiet zu Mozart geh"orende Referenzen.

Wenn man sich also dieses Wissen der Authoren zunutze macht, kann man ein Themengebiet wesentlich umfangreicher beschreiben als dies mit einfachen Textsuchmaschinen m"oglich ist.

Es bieten sich hier verschiedene M"oglichkeiten der Linkauswertung an:
\begin{description}
	\item[direkt:]	Viele Links werden benannt nach dem Thema, auf welches sie verweisen, z.B. "`Hier finden sie etwas "uber \underline{Mozart}"'. 
	
	Wenn ich nun nach "`Mozart"' suche, ben"otige ich nur alle Referenzen in denen das Wort "`Mozart"' auftaucht. Zum besseren Verst"andnis sei auf Abbildung \ref{LinksErsteSeite} verwiesen.

\begin{figure}[p]
	\makebox[\textwidth]{\epsfxsize=\textwidth \epsffile{LinksErsteSeite.eps} }
	\caption{Direkte Auswertung der Namen von Verweisen.}
	\label{LinksErsteSeite}
\end{figure}

	Wenn viele Dokumente auf ein Dokument verweisen, kann man einfach feststellen, ob die Seite, auf die verwiesen wird (in der Abbildung das Dokument $D$), eine wichtige Seite ist. 

\begin{figure}[p]
	\makebox[\textwidth]{\epsfxsize=\textwidth \epsffile{LinksZweiteSeite.eps} }
	\caption{Indirekte Auswertung von Verweisen.}
	\label{LinksZweiteSeite}
\end{figure}

	\item[indirekt:]	Alle Hypertextdokumente, die auf ein Hypertextdokument verweisen, haben zum Gro"steil auch etwas mit dessen Inhalt zu tun. So wird man auf einer Seite "uber Mozart selten einen Verweis auf {\em Gummistiefel } finden. Ein Bild zu diesem Konzept liefert Abbildung \ref{LinksZweiteSeite}.
	
	Schwieriger wird es schon, die Informationen auszuwerten, die einem die Seiten liefern, auf die bei diesem Verfahren gezeigt wird. Eine Auswertung von W"ortern die auf allen so verwiesenen Seiten zu finden waren lieferte aber schon ansprechende Ergebnisse.		
\end{description}


\section{Link-Cluster}
Die Benutzung der Links als Informationstr"ager und das Normalisieren der Dokumente bringt schon eine deutlich verbesserte Qualit"at der Suchantworten.

Die verbesserte Qualit"at ist darauf zu"urckzuf"uhren, da"s Dokumente, die nicht mit dem Suchbegriff schon von anderen Dokumenten referenziert wurden, nicht auftauchen.

Um eine weitere Qualit"at den Suchantworten hinzuzuf"ugen, entschied ich mich die normalisierten Links der Zieldokumente bei Suchantworten mit auszugeben. An einem Beispiel wird dies einfacher ersichtlich:
\begin{center}
\begin{tabular}{lr}
DokumentA enth"alt:	&"`Mozart"' auf DokumentB	\\
					&"`Mozart"' auf DokumentC	\\
DokumentB enth"alt:	&"`Klassische Komponisten"'auf DokumentA	\\
					&"`Beethoven"' auf DokumentC	\\
DokumentC enth"alt:	&"`Mozart"'auf DokumentB	\\
					&"`Klassik"' auf DokumentA	\\
\end{tabular}
\end{center}
Bei einer Suchanfrage nach "`Mozart"' erhielt man daraufhin:
\begin{center}
\begin{tabular}{lr}
Mozart:				&DokumentB	\\
					&DokumentC	\\
Komponisten:		&DokumentA	\\
Beethoven:			&DokumentC	\\
Klassik				&DokumentA	\\
\end{tabular}
\end{center}

Die erreichte Qualit"at ist damit schon eine ganz andere, man sieht nun auch "`nahe Verwandte"' des Themenkomplexes.


\section{Frequenzlisten}

Bei einem weiteren Ausbaus des "`Link-Cluster"'-Ansatzes benutzte ich die in den gefundenen Dokumenten vorkommenden W"orter als Gesamtw"orterbuch. Aus diesem W"orterbuch generierte ich dann eine Frequenzliste.

Hintergedanke dabei war, da"s Themengebiete Schlagw"orter haben, die h"aufiger vorkommen als alle anderen W"orter. Die Generierung einer Frequenzliste ist mit wesentlich weniger Aufwand verbunden, wie die Clusterung der Daten.

Bei einer Anfrage nach "`Bundeskanzler"' tauchten die folgenden W"orter dann hochfrequent auf: Saarland, Lafontaine,
Bundestag, Kohl,\dots .

Interessant ist diese Ergebnis insofern, als da"s der Name Lafontaine h"aufiger auftauchte als der von Kohl. Dies kann aber auch durch die sehr geringe Datenmenge (200 Megabyte Text von deutschen Web-Sites) , die dieser Analyse zugrunde lag, hervorgerufen worden sein.


\section{Ausblick}

Erste Versuche mit den Verweisauswertungen und dessen Verbesserungen zeigten aufregende Ergebnisse. So wurde bei der Anfrage nach "`Bundesliga"', ganz wider Erwarten, neben den "`normalen"' Ergebnisseiten auch auf die Wirtschaftsuniversit"at von Frankfurt am Main verwiesen, da diese den gr"o"sten Fu"sballergebnisserver betreibt.

Anhand solcher Ergebnisse wird klar, da"s eine Linkauswertung mehr der Denkstruktur von Menschen "ahnelt. Viele Menschen werden, zum Thema Bundesliga befragt,die Vereine und deren Spieler nennen, nur wenige werden von Datenbanken berichten.

Es l"a"st sich klar eine neue Qualit"at der Suchantworten erkennen, weshalb in einem weiteren Projekt am {\em CIS} die Verarbeitung der Linkinformationen umfangreicher untersucht wird.




